\section*{第 \ref{ch:b1:taylor} 章——泰勒公式与
渐近展开}

\begin{solution}{exo:b1:taylor:1}
\[
\eu^{2x} = 1 + 2x + 2x^2 + \frac{4x^3}{3} + o(x^3);
\qquad
\ln(1 - x) = -x - \frac{x^2}{2} - \frac{x^3}{3} - \frac{x^4}{4} +
o(x^4);
\]
\[
\sqrt{1+x} = 1 + \frac x2 - \frac{x^2}{8} + \frac{x^3}{16} + o(x^3);
\qquad
\frac{1}{1 + x^2} = 1 - x^2 + x^4 - x^6 + o(x^6),
\]
最后一式将 $-x^2$ 代入等比级数展开即得。
\end{solution}

\begin{solution}{exo:b1:taylor:2}
$\eu^x - 1 - x = \frac{x^2}{2} + o(x^2)$：极限为 $\dfrac12$。

$\cos x = 1 - \frac{x^2}{2} + \frac{x^4}{24} + o(x^4)$ 与
$\sqrt{1 - x^2} = 1 - \frac{x^2}{2} - \frac{x^4}{8} + o(x^4)$：
作差得 $\frac{x^4}{24} + \frac{x^4}{8} = \frac{x^4}{6} + o(x^4)$：
极限为 $\dfrac16$。

$\ln(1+x) - \sin x = \bigl(x - \frac{x^2}{2}\bigr) - x + o(x^2) =
-\frac{x^2}{2} + o(x^2)$：极限为 $-\dfrac12$。
\end{solution}

\begin{solution}{exo:b1:taylor:3}
$\frac{1}{1+t^2} = 1 - t^2 + t^4 + o(t^5)$；逐项\omterm{thm:b1:integration:def}{积分}，自 $0$ 至
$x$（\cref{met:b1:taylor:compute}~(5)）：
\[
\arctan x = x - \frac{x^3}{3} + \frac{x^5}{5} + o(x^5)\ \ (\text{事实上为
}o(x^6)\text{，由奇性可得}).
\]
$(1 - t^2)^{-1/2} = 1 + \frac{t^2}{2} + \frac38 t^4 + o(t^5)$（由二项式展开，
取 $\alpha = -\frac12$，$x = -t^2$：
$\binom{-1/2}{2} = \frac{(-\frac12)(-\frac32)}{2} = \frac38$）；
逐项\omterm{thm:b1:integration:def}{积分}：
\[
\arcsin x = x + \frac{x^3}{6} + \frac{3x^5}{40} + o(x^5) .
\]
\end{solution}

\begin{solution}{exo:b1:taylor:4}
由泰勒–拉格朗日余项（\cref{thm:b1:taylor:lagrange}），对 $\exp$ 在 $a =
0$，$x = 1$：$(n+1)$ 阶导数为 $\eu^t \leq \eu < 3$，在
$\intcc{0}{1}$ 上，故
\[
\Bigl|\eu - \sum_{k=0}^{n} \frac{1}{k!}\Bigr| \leq
\frac{3}{(n+1)!} .
\]
要使准确到 $6$ 位小数，需 $\frac{3}{(n+1)!} < 5\times 10^{-7}$，
即 $(n+1)! > 6\times 10^{6}$：由于 $10! = 3\,628\,800$ 及 $11! =
39\,916\,800$，故 $n + 1 = 11$，即 $n = 10$ 即可。
\end{solution}

\begin{solution}{exo:b1:taylor:5}
$n\ln\bigl(1 + \frac1n\bigr) = n\Bigl(\frac1n - \frac{1}{2n^2} +
\frac{1}{3n^3} + o\bigl(\frac{1}{n^3}\bigr)\Bigr) = 1 - \frac{1}{2n}
+ \frac{1}{3n^2} + o\bigl(\frac{1}{n^2}\bigr)$。取指数，令
$u = -\frac{1}{2n} + \frac{1}{3n^2}$，且 $\eu^u = 1 + u + \frac{u^2}2
+ o(u^2)$：
\[
\Bigl(1 + \frac1n\Bigr)^n = \eu\cdot \eu^{u}
= \eu\Bigl(1 - \frac{1}{2n} + \frac{1}{3n^2} + \frac{1}{8n^2}
+ o\Bigl(\frac{1}{n^2}\Bigr)\Bigr)
= \eu\Bigl(1 - \frac{1}{2n} + \frac{11}{24n^2} +
o\Bigl(\frac{1}{n^2}\Bigr)\Bigr).
\]
极限为 $\eu$；误差为 $\sim \dfrac{\eu}{2n}$：收敛很慢（$n$ 每增大
十倍，仅多得一位数字）。
\end{solution}

\begin{solution}{exo:b1:taylor:6}
$f(x) = x^2 - x^4 = x^2(1 + o(1))$：首项为 $x^2$，$p = 2$ 为偶数，
系数 $> 0$，故在 $0$ 处取局部极小值（并非整体：$f(2) = -12$）。

$g(x) = x^3 + x^5$：首项为 $x^3$，$p$ 为奇数，无极值；$g$
穿过其（水平）切线，$0$ 为拐点。

$h = \exp$ 在 $a = 1$ 处：$h(x) = \eu + \eu(x-1) + \frac{\eu}{2}(x-1)^2
+ o((x-1)^2)$；它与切线之差为
$\frac{\eu}{2}(x-1)^2 + o(\cdot) > 0$，局部成立。整体而言：$\eu^x -
\eu x \geq 0$ 对一切 $x$ 成立，由凸性
（\cref{thm:b1:derivative:convexchar}~(3)），图像位于每条
切线之上，仅在切点处取等号。
\end{solution}

\begin{solution}{exo:b1:taylor:7}
当 $x \to +\infty$ 时：
\[
f(x) = x\Bigl(1 + \frac1x\Bigr)^{1/3}
= x\Bigl(1 + \frac{1}{3x} - \frac{1}{9x^2} +
o\Bigl(\frac{1}{x^2}\Bigr)\Bigr)
= x + \frac13 - \frac{1}{9x} + o\Bigl(\frac1x\Bigr):
\]
渐近线为 $y = x + \frac13$，曲线在 $+\infty$ 附近位于其下方。当 $x
\to -\infty$ 时，同样的计算仍成立（立方根
对一切实数有定义，且 $\frac1x \to 0$）：渐近线仍为 $y = x +
\frac13$，但此时 $-\frac{1}{9x} > 0$：曲线\emph{位于直线上方}。
\end{solution}

\begin{solution}{exo:b1:taylor:8}
$u_n = \sqrt n\bigl(\sqrt{1 + \tfrac1n} - 1\bigr) = \sqrt
n\bigl(\frac{1}{2n} + o(\frac1n)\bigr) \sim \dfrac{1}{2\sqrt n}$.

$v_n = \ln\bigl(1 + \frac1n\bigr) \sim \dfrac 1n$.

$w_n$：令 $h = \frac1n \to 0$，$\sin h - \tan h = \bigl(h -
\frac{h^3}{6}\bigr) - \bigl(h + \frac{h^3}{3}\bigr) + o(h^3) =
-\frac{h^3}{2} + o(h^3)$，故 $w_n \sim -\dfrac{1}{2n^3}$。
\end{solution}

\begin{solution}{exo:b1:taylor:9}
用 $1$ 阶泰勒–拉格朗日公式在 $x$ 的两侧展开：
\[
f(x + h) = f(x) + h f'(x) + R_+,\quad
f(x - h) = f(x) - h f'(x) + R_-,
\qquad \abs{R_\pm} \leq \frac{h^2}{2} \sup \abs{f''} .
\]
两式相减：$f(x+h) - f(x-h) = 2h f'(x) + (R_+ - R_-)$，于是
\[
\abs{f'(x)} \leq \frac{\abs{f(x+h) - f(x-h)}}{2h} +
\frac{h}{2}\sup_{\intcc{x-h}{x+h}}\abs{f''} .
\]
由全局界可得：$\abs{f'(x)} \leq \frac{M_0}{h} + \frac{M_2
h}{2}$ 对每个 $h > 0$ 成立。右端在 $h =
\sqrt{2M_0/M_2}$（导数为零），此时取值为 $\sqrt{2M_0M_2}$
——因此 $\abs{f'} \leq \sqrt{2M_0M_2}$。（若 $M_2 = 0$，令 $h \to
\infty$：$f' = 0$，与此相符。）
\end{solution}

\begin{solution}{exo:b1:taylor:10}
在 $\intoo{0}{\pi}$ 上：$0 < \sin u < u$，故 $(u_n)$ 严格
递减、为正，因而收敛；其极限是一个
$\sin$ 在 $\intcc{0}{\pi}$ 中的不动点，由 $\sin \ell = \ell$ 推出 $\ell =
0$（因 $\sin x < x$ 对 $x > 0$ 成立）。
\begin{enumerate}
  \item 利用 $\sin u = u - \frac{u^3}{6} + o(u^3)$，当 $u \to 0$ 时：
        \[
        v_{n+1} - v_n = \frac{1}{\sin^2 u_n} - \frac{1}{u_n^2}
        = \frac{1}{u_n^2}\Bigl(\Bigl(1 - \frac{u_n^2}{6} +
        o(u_n^2)\Bigr)^{\!-2} - 1\Bigr)
        = \frac{1}{u_n^2}\Bigl(\frac{u_n^2}{3} + o(u_n^2)\Bigr)
        \longrightarrow \frac13 .
        \]
  \item 由\cref{exo:b1:seq:10}~(3)（差分形式的切萨罗定理）得，
        $\frac{v_n}{n} \to \frac13$，即 $v_n \sim \frac n3$，亦即
        $u_n^2 \sim \frac 3n$：由 $u_n > 0$，
        \[
        u_n \sim \sqrt{\frac{3}{n}} .
        \]
\end{enumerate}
\end{solution}

\begin{solution}{exo:b1:taylor:11}
通分。第一个极限：
\[
\frac{1}{x^2} - \frac{1}{\sin^2 x}
= \frac{\sin^2 x - x^2}{x^2\sin^2 x},
\qquad
\sin^2 x = \Bigl(x - \frac{x^3}{6} + o(x^4)\Bigr)^{\!2}
= x^2 - \frac{x^4}{3} + o(x^5) ,
\]
分子为 $-\frac{x^4}{3} + o(x^4)$，而
分母 $\sim x^4$，故极限为 $-\dfrac13$。

第二个式子：$\dfrac1x - \dfrac{1}{\ln(1+x)} = \dfrac{\ln(1+x) -
x}{x\ln(1+x)}$；分子为 $-\frac{x^2}{2} + o(x^2)$，
分母 $x\bigl(x + o(x)\bigr) \sim x^2$，故极限为
$-\dfrac12$。
\end{solution}

\begin{solution}{exo:b1:taylor:12}
在 $I_k = \intoo{k\pi - \frac\pi2}{k\pi + \frac\pi2}$ 上，
函数 $g(x) = \tan x - x$ 的导数为 $\tan^2 x \geq 0$，
且仅在单点 $k\pi$ 处为零，故 $g$ 在 $I_k$ 上严格
递增（由\cref{cor:b1:derivative:monotone}~(2)），
两端极限分别为 $-\infty$ 与 $+\infty$，故恰有一个
零点 $x_k$。当 $k \geq 1$ 时，$g(k\pi) = -k\pi < 0$，故 $x_k \in
\intoo{k\pi}{k\pi + \frac\pi2}$：记 $x_k = k\pi + \frac\pi2 -
\varepsilon_k$，其中 $\varepsilon_k \in \intoo{0}{\frac\pi2}$。
于是
\[
x_k = \tan x_k = \tan\Bigl(\frac\pi2 - \varepsilon_k\Bigr)
= \frac{1}{\tan\varepsilon_k}
\quad\Longrightarrow\quad
\varepsilon_k = \arctan\frac{1}{x_k} ,
\]
利用 $\tan\varepsilon_k = \frac{1}{x_k}$ 及 $\varepsilon_k \in
\intoo{0}{\frac\pi2}$。由于 $x_k \geq k\pi \to \infty$，
$\varepsilon_k \to 0$，且
\[
\varepsilon_k = \arctan\frac{1}{x_k} = \frac{1}{x_k} +
O\Bigl(\frac{1}{x_k^3}\Bigr)
= \frac{1}{k\pi + O(1)} + O\Bigl(\frac{1}{k^3}\Bigr)
= \frac{1}{k\pi} + O\Bigl(\frac{1}{k^2}\Bigr) ,
\]
由此 $x_k = k\pi + \frac\pi2 - \frac{1}{k\pi} +
o\bigl(\frac1k\bigr)$。
\end{solution}

\begin{solution}{pb:b1:taylor:1}
\textbf{1.} $S_{2n+1} - S_{2n-1} = a_{2n} - a_{2n+1} \geq 0$，且
$S_{2n+2} - S_{2n} = a_{2n+2} - a_{2n+1} \leq 0$，而 $S_{2n}
- S_{2n+1} = a_{2n+1} \to 0$：数列 $(S_{2n+1})$、
$(S_{2n})$ 为\omterm{thm:b1:seq:adjacent}{相邻数列}，收敛于公共极限 $S$
（见\cref{thm:b1:seq:adjacent}），且 $S_{2n+1} \leq S \leq S_{2n}$。
当 $n$ 为偶数时，$S_{n+1} \leq S \leq S_n$ 给出 $-a_{n+1} \leq S -
S_n \leq 0$；当 $n$ 为奇数时，$0 \leq S - S_n \leq a_{n+1}$。两种情形下，
$\abs{S - S_n} \leq a_{n+1}$，且 $S - S_n$ 的符号与
$(-1)^{n+1}$ 相同，即被略去的第一项的符号。严格递减使
所列各不等式均严格成立，特别地 $0 < \abs{S -
S_n} < a_{n+1}$。

\textbf{2.} $a_k = \frac{1}{k!}$ 严格递减且趋于 $0$：
第 1 题的结论适用。利用泰勒–拉格朗日公式
（见\cref{thm:b1:taylor:lagrange}），对 $\exp$ 在 $-1$ 与 $0$ 之间：
$\abs{\eu^{-1} - T_n} \leq \frac{1}{(n+1)!}$（导数
$\eu^t$ 在此处 $\leq 1$），故 $T_n \to \eu^{-1}$，而
第 1 题的极限 $S$ \emph{即是} $\eu^{-1}$，且具有严格界
$0 < \abs{\eu^{-1} - T_n} < \frac{1}{(n+1)!}$。

\textbf{3.} 将恒等式在 $\intcc{0}{1}$ 上\omterm{thm:b1:integration:def}{积分}：
左边为 $\arctan 1 = \frac\pi4$（由微\omterm{thm:b1:integration:def}{积分}基本定理），
第 $k$ 项给出 $\frac{(-1)^k}{2k+1}$，于是
\[
\rho_n = (-1)^{n+1}\int_0^1 \frac{t^{2n+2}}{1+t^2}\dd t,
\qquad
\abs{\rho_n} \leq \int_0^1 t^{2n+2}\dd t = \frac{1}{2n+3} .
\]

\textbf{4.} $\pi$ 的误差为 $4\abs{\rho_n} \leq
\frac{4}{2n+3}$：要使其小于 $10^{-6}$，需 $2n + 3 > 4\cdot10^6$，
约需二百万项。同时 $4S_4 = 4\bigl(1 - \frac13 +
\frac15 - \frac17 + \frac19\bigr) = 4 \times 0.834921 =
3.339683$，相差近 $0.2$，离 $\pi$ 尚远：五项，甚至
连一位数字都不能确定。

\textbf{5.} 在 $\intcc{0}{1}$ 上\omterm{thm:b1:integration:def}{积分}：$\ln 2 =
\sum_{k=1}^{n} \frac{(-1)^{k-1}}{k} + (-1)^n R_n$，其中 $R_n =
\int_0^1 \frac{t^n}{1+t}\dd t$；由 $\frac12 \leq \frac{1}{1+t}
\leq 1$ 可得 $\frac{1}{2(n+1)} \leq R_n \leq \frac{1}{n+1}$。
误差被夹在两个 $\frac1n$ 的倍数之间：六位
小数约需一百万项。

\textbf{6.} 自 $0$ 到 $x$ \omterm{thm:b1:integration:def}{积分}：$\frac12\ln\frac{1 +
x}{1 - x} = \sum_{k=0}^{n} \frac{x^{2k+1}}{2k+1} + \int_0^x
\frac{t^{2n+2}}{1-t^2}\dd t$。当 $x = \frac13$ 时：$\frac{1 +
1/3}{1 - 1/3} = 2$，且在 $\intcc{0}{\frac13}$ 上，$\frac{1}{1 -
t^2} \leq \frac98$：
\[
\ln 2 = 2\sum_{k=0}^{n} \frac{(1/3)^{2k+1}}{2k+1} +
\tilde\rho_n,
\qquad
0 < \tilde\rho_n \leq \frac{9}{4}\cdot
\frac{(1/3)^{2n+3}}{2n+3} :
\]
每多取一项，误差约除以 $9$。

\textbf{7.} 归纳法：当 $n = 1$ 时，$1 - \frac12 = \frac12 = H_2
- H_1$。归纳步：
\[
H_{2n+2} - H_{n+1} = (H_{2n} - H_n) + \frac{1}{2n+1} +
\frac{1}{2n+2} - \frac{1}{n+1}
= (H_{2n} - H_n) + \frac{1}{2n+1} - \frac{1}{2n+2},
\]
这恰好是交错和的增量。
$H_{2n} - H_n = \sum_{k=1}^{n}\frac{1}{n+k}$ 是
\cref{ex:b1:integration:riemannexample}中的\omterm{thm:b1:integration:riemann}{黎曼和}，收敛于 $\ln
2$：方法一的偶数部分和\emph{即是}
方法三的\omterm{thm:b1:integration:riemann}{黎曼和}。

\textbf{8.} $\sum_{k=1}^{6}\frac{(-1)^{k-1}}{k} = 0.61667$，
误差 $0.0765$；方法二在 $n = 5$ 时给出 $0.6931471$，误差
$\leq \frac94\cdot\frac{(1/3)^{13}}{13} = 1.1\cdot10^{-7}$。
原因在于：方法一在\omterm{def:b1:topology:closure}{边界}
点 $x = 1$ 处求对数级数的值，此处各项仅按 $\frac1k$ 的速度衰减；方法二
在 $x = \frac13$ 处求值，远离\omterm{def:b1:topology:closure}{边界}，此处每一项都带有一个新的
因子 $\frac19$。

\textbf{9.} 要求十位小数，即需 $\tilde\rho_n \leq
5\cdot10^{-11}$。当 $n = 10$ 时：$\frac94 \cdot
\frac{(1/3)^{23}}{23} = \frac94\cdot\frac{1.06\cdot10^{-11}}{23}
\approx 1.0\cdot10^{-12} < 5\cdot10^{-11}$，十一项
即已足够。

\textbf{10.} $(5+\iu)^2 = 24 + 10\iu$，于是 $(5+\iu)^4 = (24 +
10\iu)^2 = 476 + 480\iu$；又 $2(1+\iu)(239+\iu) = 2(238 +
240\iu) = 476 + 480\iu$，二者相等。辐角：$\arg(5 + \iu) =
\arctan\frac15$，故左边的辐角为
$4\arctan\frac15 \approx 0.79 \in \intoo{0}{\pi}$；右边
的辐角为 $\frac\pi4 + \arctan\frac{1}{239} \in
\intoo{0}{\pi}$。两个相等的复数，其辐角都位于
长度 $< 2\pi$ 的同一\omterm{prop:b1:reals:intervals}{区间}内：
\[
4\arctan\frac15 = \frac\pi4 + \arctan\frac{1}{239} ,
\]
这就是马青公式。

\textbf{11.} 将 $\frac{1}{1+t^2} = \sum_{k=0}^n (-1)^k
t^{2k} + \frac{(-1)^{n+1}t^{2n+2}}{1+t^2}$ 自 $0$ 到 $x$ \omterm{thm:b1:integration:def}{积分}：
\[
\arctan x = \sum_{k=0}^{n}\frac{(-1)^k x^{2k+1}}{2k+1} +
r_n(x),
\qquad
\abs{r_n(x)} \leq \int_0^x t^{2n+2}\dd t =
\frac{x^{2n+3}}{2n+3} .
\]

\textbf{12.} 误差为：$16\,\frac{(1/5)^{11}}{11} = 3.0\cdot
10^{-8}$ 与 $4\,\frac{(1/239)^5}{5} \approx 1\cdot10^{-12}$：
合计 $< 5\cdot10^{-8}$。上式所示之和的值为
$3.14159268\dots$，因此 $\pi = 3.1415926\dots$ 可确认精确到
$5\cdot10^{-8}$：仅用六项即得七位小数（其中
五项取 $\frac15$，两项取 $\frac1{239}$，宽松计算）。

\textbf{13.} 莱布尼茨：误差 $\sim \frac1n$，因此每多一位数字，
工作量便增至十倍。达尔泽尔
（见\cref{pb:b1:integration:1}中第 22 题）：误差 $4^{1-5m}$，
每步约增加三位数字，但每步的\omterm{def:b1:poly:def}{多项式}都更繁重。
马青：每项的误差比为 $\frac{1}{25}$，约得 $1.4$ 位数字，
每项只需一次除法。其中的道理是：几何型展开的余项
的量级为 $x^{2n}$，故使 $x$
足够小，便可在每项\emph{固定}的代价下换得更多数字——马青的
复数恒等式正是一台缩小 $x$ 的机器。

\textbf{14.} $a_k = \frac{1}{(2k)!}$ 严格递减趋于 $0$；
由第 1 题及\cref{prop:b1:taylor:standard}（拉格朗日
余项估计，同第 2 题），$\sum_{k \leq n}\frac{(-1)^k}{(2k)!}
\to \cos 1$，并有\emph{严格}的夹逼式 $0 < \bigl|\cos 1 -
S'_n\bigr| < \frac{1}{(2n+2)!}$。设 $\cos 1 = \frac pq$，
并取 $2n \geq q$，则 $(2n)!\,S'_n = \sum_{k\leq n} (-1)^k
\frac{(2n)!}{(2k)!} \in \Z$ 且 $(2n)!\,\frac pq \in \Z$，而
\[
0 < \Bigl|(2n)!\cos 1 - (2n)!S'_n\Bigr| <
\frac{(2n)!}{(2n+2)!} = \frac{1}{(2n+1)(2n+2)} < 1 :
\]
绝对值 $< 1$ 的一个非零整数。矛盾，故 $\cos
1 \notin \Q$。

\textbf{15.} 类似地，取 $a_k = \frac{1}{(2k+1)!}$，
乘以 $(2n+1)!$，取 $2n + 1 \geq q$，即得 $\sin 1 \notin
\Q$。然而 $\cos^2 1 + \sin^2 1 = 1 \in \Q$：无理数的和与积
可以是有理数——无理性不会
免费地穿过任何代数运算。

\textbf{16.} 尾项估计：对 $m > n$，
\[
\sum_{k=n+1}^{m} \frac{1}{(2k)!}
\leq \frac{1}{(2n+2)!}\Bigl(1 + \frac12 + \frac14 +
\dots\Bigr) \leq \frac{2}{(2n+2)!} ,
\]
因为各相邻两项之比为 $\frac{1}{(2k+1)(2k+2)} \leq
\frac12$；尾项为正（其首项即为正）。于是 $0 <
\cosh 1 - \sum_{k\leq n}\frac{1}{(2k)!} < \frac{2}{(2n+2)!}$,
再乘以 $(2n)!$，取 $2n \geq q$，便把一个非零
整数再次夹入 $\intoo{0}{1}$ 中，故 $\cosh 1 \notin \Q$。

\textbf{17.} $\cos\frac1m = \sum_k
\frac{(-1)^k}{m^{2k}(2k)!}$：各项严格递减，趋于
零，且 $m^{2n}(2n)!\cdot\frac{1}{m^{2k}(2k)!} =
m^{2(n-k)}\frac{(2n)!}{(2k)!} \in \Z$，对 $k \leq n$。
若 $\cos\frac1m = \frac pq$，则将严格夹逼式乘以
$q\,m^{2n}(2n)!$：误差
以 $\frac{q}{m^2(2n+1)(2n+2)} < 1$ 为界（$n$ 充分大时），矛盾。
对 $\frac ab$（$a \geq 2$），去分母即将
尾项乘以 $b^{2n}(2n)!$，但被略去的
首项为 $\frac{a^{2n+2}}{b^{2n+2}(2n+2)!}$，两者之积
$\frac{a^{2n+2}}{b^2(2n+1)(2n+2)}$ \emph{剧增}：
分子 $a^{2n+2}$ 不再能消去分母，夹逼随之失效。
（结论仍然成立——借助尼文式的工具，而非
此法。）

\textbf{18.} 当 $t = x$ 时，$g$ 的每一项除 $f(x)
- f(x) = 0$ 外均为零，故 $g(x) = 0$；$A$ 的取法使 $g(a) = 0$。
求导后，和式裂项相消：
\[
g'(t) = -\frac{f^{(n+1)}(t)}{n!}(x - t)^n +
A\,\frac{(x-t)^n}{n!}
= \frac{(x-t)^n}{n!}\bigl(A - f^{(n+1)}(t)\bigr) .
\]
在 $a$ 与 $x$ 之间的\omterm{prop:b1:reals:intervals}{区间}上用罗尔定理，得 $c$ 严格介于其间，
使 $g'(c) = 0$；由 $(x - c)^n \neq 0$ 得 $A =
f^{(n+1)}(c)$。展开 $g(a) = 0$ 即得泰勒等式，
其余项为 $\frac{f^{(n+1)}(c)}{(n+1)!}(x-a)^{n+1}$。

\textbf{19.} 当 $x > 0$ 时，余项为 $\frac{\eu^{c}}
{(n+1)!}x^{n+1} > 0$：指数函数严格超过其每个
泰勒\omterm{def:b1:poly:def}{多项式}，每一阶均如此。当 $x < 0$ 时，
余项的符号与 $x^{n+1}$ 相同：$\eu^x$ 在
\emph{高于}其\omterm{def:b1:poly:def}{多项式}（$n$ 为奇数时），\emph{低于}之
（$n$ 为偶数时）——两侧交替，正如 $1 + x$ 与 $1
+ x + \frac{x^2}{2}$ 相对于 $\eu^x$ 的图像已经显示。

\textbf{20.} 真实误差：$\sin 0.5 - 0.4791667 =
2.59\cdot10^{-4}$，与界 $\frac{0.5^5}{120} =
2.60\cdot10^{-4}$：几乎取到（后一项主导了
尾项）。杨氏型：用于极限与局部分析，无需任何
常数。拉格朗日型：用于有保证的小数。交错型：
可用时，界与拉格朗日型相同，并\emph{另外}给出误差的方向——
三者中最优，但最少见。

\textbf{21.} 在 $0$ 处\omterm{def:b1:continuity:continuous}{连续}：令 $u = \frac{1}{x^2} \to
+\infty$，$f(x) = \eu^{-u} \to 0 = f(0)$。$0$ 处的导数：
$\bigl|\frac{f(h)}{h}\bigr| = \sqrt u\,\eu^{-u} \to 0$
（见\cref{prop:b1:functions:powerrules}）：$f'(0) = 0$。当 $x \neq
0$ 时，$f'(x) = \frac{2}{x^3}\eu^{-1/x^2}$：形如
$P_1\bigl(\frac1x\bigr)\eu^{-1/x^2}$，其中 $P_1(X) = 2X^3$；
由归纳，对 $P_k(\frac1x)\eu^{-1/x^2}$ 求导得
$P_{k+1}(X) = 2X^3 P_k(X) - X^2 P_k'(X)$，为\omterm{def:b1:poly:def}{多项式}。于是
\[
\frac{f^{(k)}(h) - 0}{h} = v\,P_k(v)\,\eu^{-v^2}
\Big|_{v = 1/h} \longrightarrow 0
\]
（\omterm{def:b1:poly:def}{多项式}与 $\eu^{-v^2}$ 在
$\pm\infty$ 处的\omterm{prop:b1:functions:powerrules}{增长比较}）。由归纳法，$f^{(k)}(0) = 0$ 对一切 $k$ 成立，于是
$f$ 在 $0$ 处的泰勒\omterm{def:b1:poly:def}{多项式}全部为零，然而 $f > 0$ 在 $0$ 之外：
泰勒--杨公式在每一阶上都精确，却对
芽之外全然无知。展开式只是局部信息。

\textbf{22.} 由第 2 题，$0 < \abs{\eu^{-1} - T_n} <
\frac{1}{(n+1)!}$，严格成立。若 $\eu^{-1} = \frac pq$，取 $n
\geq q$，并乘以 $n!$：$n!\,T_n \in \Z$ 且 $n!\frac pq
\in \Z$，于是一个非零整数的绝对值 $<
\frac{n!}{(n+1)!} = \frac{1}{n+1} < 1$，矛盾。因此
$\eu^{-1} \notin \Q$，从而 $\eu = \frac{1}{\eu^{-1}}$ 亦为
无理数。三种证明：以\omterm{thm:b1:seq:adjacent}{相邻数列}
夹逼 $q!\,\eu$（见\cref{exo:b1:seq:9}）；\omterm{thm:b1:integration:def}{积分}递推 $A_n =
\eu - nA_{n-1}$（见\cref{pb:b1:integration:1}）；交错
夹逼（此处）。一个陷阱，三种刻画“小量”的凭据。

\textbf{23.} 将部分和两两分组：$\sum_{k=1}^{2n}
(-1)^k b_k = E_n - O_n$，其中 $E_n = \sum_{j=1}^{n}
\frac{1}{4j^2}$，有界（由
\cref{ex:b1:seq:basel}的裂项估计，$E_n \leq \frac12$），且 $O_n =
\sum_{j=1}^{n}\frac{1}{2j-1} \geq \frac12 H_n \to +\infty$
（见\cref{exo:b1:seq:5}）：部分和趋于 $-\infty$。
单调性恰在第 1 题中被用于 $S_{2n+1} -
S_{2n-1} = a_{2n} - a_{2n+1}$ 需定号之处：若无递减性，
偶数项与奇数项所成的\omterm{def:b1:seq:subsequence}{子列}未必单调，
相邻性便不复成立。

\textbf{24.} $(2+\iu)(3+\iu) = 5 + 5\iu = 5(1+\iu)$；两边取
辐角（均在 $\intoo{0}{\frac\pi2}$ 内）：$\arctan\frac12 +
\arctan\frac13 = \frac\pi4$。取六位小数时的级数代价：
误差 $\leq 4\bigl(\frac{(1/2)^{2n+3}}{2n+3} +
\frac{(1/3)^{2n+3}}{2n+3}\bigr)$，当 $n = 10$ 时约为
$\approx 2\cdot10^{-8} < 5\cdot10^{-7}$，共需十一项。比较：
较莱布尼茨级数优出指数量级，逊于马青
（其主导点 $\frac15$ 小于 $\frac12$）：
每项约得 $0.6$ 位，马青公式则为 $1.4$ 位。

\textbf{25.} (i) 对递减且 $a_k \to 0$ 者，交错
部分和收敛，$\abs{S - S_n} \leq a_{n+1}$，且
误差的符号与被略去的首项一致。(ii) 有限
恒等式若带显式余项，可在选定点处\emph{求值并
估界}，而极限命题仅
保证最终的接近——确证所需的是前者。
(iii) 汇总：借马青公式得 $\pi$ 精确到 $5\cdot10^{-8}$，$\ln 2$ 精确到
十位小数（借 $\frac13$ 级数），$\cos 1$、
$\sin 1$、$\cosh 1$、$\cos\frac1m$ 与 $\eu^{-1}$ 的无理性，
泰勒–拉格朗日等式，以及平坦函数的警示。(iv)
在\cref{ch:b1:series}中，第 1 题推广为交错级数判别法，
区分绝对收敛与条件收敛，
其周末问题为重排的戏剧搭好舞台，
其中途径 1 的交错调和级数是
明星见证。
\end{solution}
